\documentclass[11pt,a4paper]{article}
\usepackage[margin=2.5cm]{geometry}
\usepackage{amsmath,amssymb,booktabs}
\usepackage[T1]{fontenc}
\usepackage[utf8]{inputenc}
\title{Contoh Laporan OptiLab}
\date{\today}
\begin{document}
\maketitle
\section*{Formulasi Matematis}
\[\begin{aligned}
\max\quad Z &= 40x_{1} + 30x_{2} \\
\text{dengan kendala:} & \\
&\quad 2x_{1} + x_{2} \leq 100 \\
&\quad x_{1} + 2x_{2} \leq 80 \\
&\quad x_{1}\geq 0,\;x_{2}\geq 0
\end{aligned}\]
\section*{Hasil Penyelesaian}
Nilai objektif optimal: $Z^* = 2200$.

\[\begin{aligned}
x_{1} &= 40 \\
x_{2} &= 20
\end{aligned}\]
\subsection*{Analisis Kendala}
\begin{center}
\begin{tabular}{rrrrl}
\toprule Kendala & LHS & RHS & Slack/Surplus & Status \\
\midrule
1 & 100 & 100 & 0 & Aktif \\
2 & 80 & 80 & 0 & Aktif \\
\bottomrule
\end{tabular}
\end{center}
\noindent Iterasi solver yang dilaporkan: 2.
\section*{Metode}
Model diselesaikan sebagai Linear Programming kontinu menggunakan SciPy \texttt{linprog} dengan metode HiGHS. Hasil berlaku untuk koefisien dan asumsi yang dimasukkan.
\end{document}
